\section{Literature}
\subsection{Growth-at-Risk framework}
The Growth-at-Risk (GaR) framework, drawing from the Value-at-Risk (VaR) concept in finance, links the probability and severity of negative economic outcomes to pre-existing financial vulnerabilities \citep{imf_gar_2019}. Essentially, GaR captures the time lag between the accumulation of systemic risks and the materialization of losses. Unlike models focusing solely on expected outcomes, GaR accounts for the entire distribution of future output growth, particularly the tail of the distribution \citep{szabo2020}. This approach is crucial as systemic risks often behave non-linearly, especially during economic downturns.

For the policymaker, it is a practical means of differentiating and quantifying potential risks before they manifest. In small open economies, for instance, external shocks propagate through the financial system more aggressively, leading to amplified negative effects \citep{esrb2021}. A case in point is Peru, where the impact of external shocks was found to be twice as large as that of internal factors, such as leverage \citep{imf_gar_2019}. While macroprudential policies cannot prevent the occurrence or magnitude of these crises, they can mitigate their impact through countercyclical measures \citep{esrb2021}.

In practice, the GaR framework has many variations. For instance, \citet{aikman2019} employ a panel quantile regression, incorporating variables like credit-to-GDP growth, asset price volatility, house price growth, current account deficit, and bank capital. \citet{adrian_2019} took a different approach by using the National Financial Conditions Index (NFCI), a composite indicator of debt, equity, and money markets. They also smooth out the estimated quantiles by interpolating between them using a skewed t-distribution. Meanwhile, \citet{imf_gar_2019} construct multiple partitions comprising of three parts: the domestic price of risk (reflects interest rate markets and financial conditions), leverage (measures credit growth and macrofinancial vulnerabilities), and external conditions (like commodity prices). According to the authors \citep{imf_gar_2019}, the partitioning technique helps extract common trends and remove idiosyncratic noise, which improves the quality of the quantile regressions.

Similarly, \citet{galan2020} applied a panel quantile regression model using factors like the FCI \citet{koop_2013}, credit-to-GDP growth, house prices, and the MPI. \citeauthor{galan2020}'s (\citeyear{galan2020}) model diverges from \citet{adrian_2019} by employing a kernel-based non-parametric method to interpolate between quantiles, allowing for greater flexibility. They also emphasizes the role of macroprudential policy, considering its effects across different phases of the financial cycle.

Compared to the other authors, \citet{szabo2020} proposes a Bayesian quantile regression model, which addresses overfitting and enhances out-of-sample forecasting. They also use “system priors” to make the model more consistent with the official GDP projections. Focusing on systemic risks, \citet{szabo2020} utilized variables such as the financial cycle indicator (equivalent to FCyI), the banking prudence indicator (BPI), the financial stress index (CISS), and the global uncertainty index (GEPU).

Despite its variations, the GaR framework shares some foundational steps across applications. The general framework can be summarized in three key steps. First, it involves selecting variables that effectively reflect financial imbalances. Next, these variables are evaluated for their impact on the distribution of future growth outcomes, with particular attention to non-linear relationships. Finally, these estimated effects are discussed or compared to understand how they contribute to overall risk assessments. These choices collectively shape the framework's adaptability and flexibility, allowing for modifications to meet evolving analytical or policy needs.

Finally, the reports by the European Systemic Risk Board \citep{esrb2018, esrb2019, esrb2021, esrb2024} cover the GaR framework quite extensively. They decompose systemic risks into financial vulnerabilities and financial stress. The framework also includes structural factors, specifically related to the financial sector, control variables, and the MPI. Using a panel quantile regression, they compare the performance of different indices for assessing the systemic risks. Lastly, they compare the systemic risks across the European countries and discuss their relevance in the current macroprudential stance.

Building on the insights from the literature, we employ both panel and localized quantile regressions. We focus on key risk indicators such as financial vulnerabilities and financial stress, supplemented by structural variables like the capital adequacy ratio (CAR) and the liquidity ratio (HQLA/Assets). Since Armenia is a small open economy, we also incorporate the macroeconomic indicator (MI) to account for the impact of external shocks. After estimating conditional quantiles, we interpolate between them using both parametric and non-parametric methods. Finally, for deriving the macroprudential stance we use the distance-to-tail metric \citep{esrb2024}, but also extend that with the work of \citet{suarez2022}.

Consistent with prior studies \citep{imf_gar_2019}, we partition our data into indices to highlight common trends and filter out idiosyncratic noise. While we also consider the influence of macroprudential policy across different phases of the financial cycle \citep{galan2020}, we deliberately focus on a narrower set of indicators than some other studies \citep{aikman2019}. Unlike certain approaches \citep{szabo2020}, our framework avoids prior assumptions or other calibration steps designed to align it with official growth projections, prioritizing instead the estimation of current systemic risks over future growth forecasts.  Finally, we adapt and extend the framework to enhance its suitability for assessing the macroprudential stance.

\subsection{Key indicators of risk}
No crisis is the same, but during a crisis, the financial system has the potential to amplify shocks. The extent of this amplification will depend on the type and magnitude of existing vulnerabilities within the system. In this paper, we assess systemic risk exposures by examining the current financial vulnerabilities and stress.

\textit{financial vulnerabilities}: The FCyI \citep{hollo2012, plasil2014} is an extension of the Basel credit-to-GDP gap approach \citep{basel2010} that also incorporates elements from mortgage markets, property prices, and the trade deficit. It serves as a leading indicator for future deteriorations in GDP growth. A rising FCyI suggests that economic agents are accumulating substantial debt. It may also describe credit bubbles where excessive borrowing drives up asset prices. Furthermore, the FCyI can reflect sovereign risk exposures, such as those witnessed during the European debt crisis of 2010.

\textit{financial stress}: In contrast, the FCI follows the methodology proposed by \citet{koop_2013}, using A factor-augmented vector autoregressive model with time-varying coefficients and stochastic volatility to aggregate the financial variables into a single index. This index primarily focuses on the interest rate markets, thus reflecting both current and future expectations of output growth. However, within the context of the GaR framework, the FCI is primarily used as a measure of financial stress. We define financial stress as “periods in which risks have [materialized], reflected in stress episodes with larger asset price volatility, interest rate spreads, and deteriorated risk sentiment” \citep{esrb2021}.

In \autoref{fig:arm_systemic_risks_aa}, we plot the FCyI, FCI, MI, and real GDP growth. The key insight from this plot is that FCyI, FCI, and MI respond differently depending on the nature of the shock. During the 2007–2009 financial crisis, MI showed the first signs of deterioration, followed by increased financial stress, and finally, a shift in the financial cycle as losses materialized. Meanwhile, in 2014-2015 when Armenia was affected by the Russian oil crisis financial stress rose first, followed by changes in MI and the financial cycle. In the cases of the Covid-19 shock and the Nagorno-Karabakh war, the output growth was the first to decline, with vulnerabilities and financial stress increasing almost immediately after. Lastly, the Russo-Ukrainian war initially posed inflationary and trade-related risks for Armenia, but overall, it was a net positive shock for the economy due to the influx of businesses and skilled migrants. Financial conditions improved, and the financial cycle continued to rise.
\begin{figure*}
	\centering 
	\includegraphics[width=1\textwidth, angle=0]{Figures/arm_systemic_risks_aa.png}	
	\caption{The financial cycle index, financial conditions index, macroeconomic indicator, and real GDP growth. The FCyI, FCI, and MI are normalized to have zero mean and a standard deviation of one. Vertical areas in red represent economic disturbances.} 
	\label{fig:arm_systemic_risks_aa}%
\end{figure*}

Throughout these crises, the primary goal of macroprudential policy is to improve the resilience of the financial system, thereby mitigating the negative impact of the crises on economic growth. However, it is important to recognize that the effectiveness of such a policy also depends on the economic conditions. The impact of the policy can vary according to the current buildup of vulnerabilities (the position in the financial cycle) and the degree of financial stress (before, during, and after a crisis). Financial vulnerabilities represent risks that manifest during periods of financial stress, as such, these two indices are also closely interconnected. To account for the relationships between these variables, we include three interaction terms in our models: (1) the product of the normalized FCyI and FCI, referred to as the Crisis Index (CI); (2) the product of MPI and the normalized FCyI, termed the Cycle-Based Policy Adjustment (CyBPA); and (3) the product between MPI and normalized FCI, the Conditions-Based Policy Adjustment (CBPA).
